\documentclass[conference]{IEEEtran}

% ── Packages ──────────────────────────────────────────────────────────────────
\usepackage{amsmath,amssymb,amsfonts}
\usepackage{graphicx}
\usepackage{booktabs}
\usepackage{multirow}
\usepackage{cite}
\usepackage{url}
\usepackage{xcolor}
\usepackage{algorithmic}
\usepackage{algorithm}
\usepackage{bm}
\usepackage[hidelinks]{hyperref}

% ── Macros ────────────────────────────────────────────────────────────────────
\newcommand{\vp}{\mathbf{p}}
\newcommand{\vv}{\mathbf{v}}
\newcommand{\vdelta}{\bm{\delta}}
\newcommand{\mat}[1]{\mathbf{#1}}
\newcommand{\norm}[1]{\left\lVert#1\right\rVert}

% ── Title / Authors ───────────────────────────────────────────────────────────
\title{PI-xLSTM: A Physics-Informed Bidirectional Extended LSTM\\
for UAV SAR Trajectory Refinement}

\author{
  \IEEEauthorblockN{[Author Names Anonymized for Review]}
  \IEEEauthorblockA{[Affiliation]}
}

\begin{document}
\maketitle

% ──────────────────────────────────────────────────────────────────────────────
\begin{abstract}
Accurate antenna phase center (APC) trajectory estimation is essential for high-resolution
unmanned aerial vehicle synthetic aperture radar (UAV-SAR) imaging.
Navigation sensors provide continuous trajectory measurements,
but their accumulated errors—comprising low-frequency drift,
mechanical vibrations, and intermittent step-change faults—directly degrade
SAR phase coherence and focusing quality.
This paper proposes PI-xLSTM, a \emph{physics-informed bidirectional extended
long short-term memory} network for UAV-SAR trajectory refinement.
The model ingests a 12-dimensional feature vector constructed from raw navigation
measurements, Doppler-derived velocity, slant-range observations from sparse
range compression module (RCM) reference targets, and navigation–RCM residuals.
A compound loss function couples a normalized mean-squared error (MSE) term
with a smoothness regularizer and a novel RCM physical consistency constraint,
whose influence is progressively activated via curriculum learning to stabilize
convergence.
To bridge the training–inference gap arising from variable numbers of
available ground reference points, a random reference-point masking augmentation
is introduced during training.
A diversity-enhanced simulation pipeline generates 5{,}000 training trajectories
spanning stratified error magnitudes, random heading rotations,
lateral offsets, and intermittent step-change faults.
Preliminary results on the held-out validation set demonstrate a trajectory RMSE
of 28.6\,mm, representing a 32.5\% error reduction relative to the raw navigation
input.
\end{abstract}

\begin{IEEEkeywords}
SAR motion compensation, trajectory refinement, xLSTM, physics-informed learning,
range compression module, curriculum learning
\end{IEEEkeywords}

% ──────────────────────────────────────────────────────────────────────────────
\section{Introduction}
\label{sec:intro}

Synthetic aperture radar (SAR) aboard unmanned aerial vehicles
has attracted growing interest in precision agriculture, disaster assessment,
and urban surveillance owing to the platform's flexibility and low operational cost
\cite{cumming2005digital}.
Unlike satellite or crewed-aircraft SAR systems,
UAV platforms exhibit pronounced positional instability caused by wind gusts,
airframe vibrations, and imperfect PID attitude control.
The SAR phase history is exquisitely sensitive to antenna phase center (APC) motion:
a position error of only $\lambda/8$ (approximately 2–4\,mm at X-band)
induces a $\pi/4$ phase perturbation that measurably degrades the impulse response
function (IRF) \cite{moreira1994airborne}.

Conventional motion compensation (MoComp) pipelines combine
inertial navigation system (INS) measurements with GPS corrections and
optionally apply autofocus algorithms such as phase gradient autofocus (PGA)
\cite{wahl1994pga}.
However, accumulated IMU drift over long synthetic apertures,
GPS multipath in cluttered environments, and the scene-dependence of autofocus
limit the achievable correction accuracy for lightweight UAV systems.

Recent advances in deep sequence modeling offer an alternative:
given sufficiently diverse training data,
a neural network can learn to infer trajectory errors directly from
observable navigation features and sparse physical measurements
without explicit aerodynamic modeling.
Several works have demonstrated LSTM and Transformer networks for
dead-reckoning and INS integration
\cite{yan2018ridi,liu2022aionet},
yet their application to SAR-specific trajectory refinement—where
physical range constraints are available but ground truth may be absent at inference
time—remains largely unexplored.

\textbf{Contributions.}
This paper makes the following contributions:
\begin{enumerate}
  \item We propose \textbf{PI-xLSTM}, the first application of the extended
        LSTM (xLSTM) architecture \cite{beck2024xlstm}
        to UAV-SAR trajectory refinement,
        using a bidirectional formulation with shared weights between
        forward and backward passes.
  \item We introduce a \textbf{RCM physical consistency loss} that enforces
        agreement between the predicted trajectory and observed
        range compression module (RCM) slant ranges at sparse ground reference
        targets, embedded within a curriculum learning schedule.
  \item We propose \textbf{random reference-point masking} as a data
        augmentation strategy to close the training–inference gap when
        fewer than three reference targets are available at inference time.
  \item We design a \textbf{diversity-enhanced simulation pipeline} that generates
        trajectories with stratified error magnitudes,
        random heading rotations ($\pm15^\circ$),
        lateral offsets ($\pm100$\,m),
        and intermittent step-change fault events,
        substantially reducing overfitting to homogeneous simulated data.
\end{enumerate}

% ──────────────────────────────────────────────────────────────────────────────
\section{System Model and Problem Formulation}
\label{sec:system}

\subsection{SAR Trajectory Error Model}

Let $\vp_{\text{true}}[k] \in \mathbb{R}^{3}$, $k = 0,\ldots,T{-}1$,
denote the true APC position at pulse index $k$ in a Cartesian frame
$(Y_\text{lat},\, X_\text{fwd},\, Z_\text{up})$.
The navigation system outputs a perturbed measurement
\begin{equation}
  \vp_{\text{raw}}[k] = \vp_{\text{true}}[k] + \vdelta[k],
  \label{eq:error_model}
\end{equation}
where $\vdelta[k]$ is the APC position error comprising four
physically motivated components:
\begin{equation}
  \vdelta[k] = \vdelta_\text{trend}[k] + \vdelta_\text{vib}[k]
               + \vdelta_\text{rnd}[k]  + \vdelta_\text{step}[k].
  \label{eq:error_decomp}
\end{equation}
Specifically, $\vdelta_\text{trend}$ is a polynomial bias–drift term,
$\vdelta_\text{vib}$ contains multi-frequency mechanical vibrations
($0.5$–$3$\,Hz),
$\vdelta_\text{rnd}$ is a spectrally shaped random walk,
and $\vdelta_\text{step}$ models intermittent step-change faults
of magnitude $\pm[20, 50]$\,mm occurring 0–5 times per trajectory.

The goal of trajectory refinement is to estimate the correction
$\hat{\vdelta}[k]$ from observable data, yielding the refined trajectory
$\hat{\vp}_{\text{fix}}[k] = \vp_{\text{raw}}[k] - \hat{\vdelta}[k]$.

\subsection{RCM Physical Constraint}

Let $\mathcal{R} = \{P_A, P_B, P_C\} \subset \mathbb{R}^{3}$
denote up to three ground reference targets with known coordinates.
The true slant range from the APC to reference target $P_r$ at pulse $k$ is
\begin{equation}
  \rho_r^{\text{true}}[k] = \norm{\vp_{\text{true}}[k] - P_r}_2.
  \label{eq:true_rcm}
\end{equation}
In a real deployment, $\rho_r^{\text{obs}}[k]$ is measured directly from
the range compression module output (the peak position of the compressed
reference echo).
The navigation–RCM residual
\begin{equation}
  \epsilon_r[k] = \norm{\vp_{\text{raw}}[k] - P_r}_2 - \rho_r^{\text{obs}}[k]
  \label{eq:residual}
\end{equation}
is an observable surrogate that is nonlinearly correlated with $\vdelta[k]$.
A predicted trajectory $\hat{\vp}_{\text{fix}}[k]$ is physically consistent
if $\norm{\hat{\vp}_{\text{fix}}[k] - P_r}_2 \approx \rho_r^{\text{obs}}[k]$
for all $r$ and $k$.

% ──────────────────────────────────────────────────────────────────────────────
\section{Proposed Method: PI-xLSTM}
\label{sec:method}

\subsection{Feature Representation}
\label{sec:features}

The input to PI-xLSTM at each pulse $k$ is a 12-dimensional feature vector
$\mathbf{f}[k] \in \mathbb{R}^{12}$ organized as:
\begin{equation}
  \mathbf{f}[k] =
  \bigl[\vp_{\text{raw}}[k];\; \mathbf{v}_{\text{raw}}[k];\;
        \bm{\rho}^{\text{obs}}[k];\; \bm{\epsilon}[k] \bigr],
  \label{eq:feature}
\end{equation}
where:
\begin{itemize}
  \item $\vp_{\text{raw}}[k] \in \mathbb{R}^{3}$: smoothed navigation trajectory;
  \item $\mathbf{v}_{\text{raw}}[k] \in \mathbb{R}^{3}$: numerically differentiated
        and smoothed velocity (Gaussian window, $\sigma{=}20$ samples);
  \item $\bm{\rho}^{\text{obs}}[k] \in \mathbb{R}^{3}$: observed RCM slant ranges
        to $P_A$, $P_B$, $P_C$ (with per-sample measurement noise $\sim\mathcal{N}(0,\sigma_\rho^2)$,
        $\sigma_\rho \in [5,30]$\,mm);
  \item $\bm{\epsilon}[k] \in \mathbb{R}^{3}$: navigation–RCM residuals (\ref{eq:residual}).
\end{itemize}
All features are normalized to zero mean and unit variance using statistics
computed over the training set.
An explicit feature mask $\mathbf{m}[k] \in \{0,1\}^{12}$ indicates which
channels are valid (zero for missing reference-point features).

\subsection{Bidirectional xLSTM Architecture}
\label{sec:arch}

The extended LSTM (xLSTM) \cite{beck2024xlstm} enriches the standard LSTM
with two complementary cell types:
the \emph{matrix LSTM} (mLSTM), which replaces the scalar cell state with
a low-rank matrix cell that enables fully parallelizable training via
quadratic attention, and the
\emph{scalar LSTM} (sLSTM), which retains the classical scalar state with
exponential gating and a hardware-optimized CUDA kernel.
The xLSTM block stack interleaves $B{-}1$ mLSTM blocks with one sLSTM block
(at position 1), which we find to provide a stabilizing inductive bias for
time-series regression.

\textbf{Bidirectional extension.}
SAR trajectory refinement benefits from bidirectional context:
errors visible in future pulses can help disambiguate the correction at the
current pulse.
We construct a bidirectional stack by running the \emph{same} weight-shared
xLSTM stack in both the forward and the time-reversed direction:
\begin{align}
  \mathbf{h}_k^{(\text{fwd})} &= \text{xLSTMStack}(\mathbf{x}_{1:T})[k], \\
  \mathbf{h}_k^{(\text{bwd})} &= \text{xLSTMStack}(\mathbf{x}_{T:1})[T{-}k{+}1],
\end{align}
where $\mathbf{x}_k = W_\text{in}\,\mathbf{f}[k]$, $W_\text{in}\in\mathbb{R}^{D\times F}$.
The outputs are concatenated and projected:
\begin{equation}
  \hat{\vdelta}[k] = W_\text{out}\,\text{Dropout}\bigl(
    [\mathbf{h}_k^{(\text{fwd})};\, \mathbf{h}_k^{(\text{bwd})}]\bigr) + b_\text{out}.
  \label{eq:output}
\end{equation}
Weight sharing reduces the parameter count by $2\times$ compared to independent
forward and backward stacks, while still capturing asymmetric temporal context.

\textbf{Hyperparameters.}
We use $B{=}4$ blocks, embedding dimension $D{=}128$, $H{=}4$ heads, and
dropout rate $p{=}0.1$, yielding approximately $1.5$\,M trainable parameters.
Input features ($F{=}12$) are projected to $D$ via a learned linear mapping.

\subsection{Physics-Informed Compound Loss}
\label{sec:loss}

The total training loss is
\begin{equation}
  \mathcal{L} = \mathcal{L}_\text{MSE}
              + \lambda_s\,\mathcal{L}_\text{smooth}
              + \lambda_\text{rcm}(e)\,\mathcal{L}_\text{RCM},
  \label{eq:total_loss}
\end{equation}
where $e$ denotes the current training epoch.

\textbf{Normalized MSE.}
Supervision is applied in the normalized error space to equalize the
per-axis loss scales:
\begin{equation}
  \mathcal{L}_\text{MSE} = \frac{1}{BT}
    \sum_{b,k}\norm*{\frac{\hat{\vdelta}_b[k] - \vdelta_b[k]}{\bm{\sigma}_\delta}}_2^2,
  \label{eq:mse}
\end{equation}
where $\bm{\sigma}_\delta$ is the per-axis error standard deviation of the training set.

\textbf{Smoothness regularizer.}
Trajectory corrections should be temporally smooth to avoid introducing
high-frequency artifacts:
\begin{equation}
  \mathcal{L}_\text{smooth} = \frac{1}{BT}
    \sum_{b,k}\norm{\hat{\vdelta}_b[k+1] - 2\hat{\vdelta}_b[k] + \hat{\vdelta}_b[k-1]}_2^2.
  \label{eq:smooth}
\end{equation}
The fixed weight $\lambda_s{=}2.0$ was selected by grid search.

\textbf{RCM physical consistency loss.}
Let $\hat{\vp}_b[k] = \vp_{\text{raw},b}[k] - \hat{\vdelta}_b[k]$ be the
predicted refined trajectory.
The RCM loss penalizes deviation of predicted slant ranges from
the observed RCM measurements:
\begin{equation}
  \mathcal{L}_\text{RCM} = \frac{1}{BT|\mathcal{R}|}\sum_{b,k,r}
    w_{b,r}\Bigl(\norm{\hat{\vp}_b[k] - P_r}_2 - \rho_r^{\text{obs}}[k]\Bigr)^2,
  \label{eq:rcm_loss}
\end{equation}
where $w_{b,r} \in \{0,1\}$ is the per-reference-point availability mask.

\textbf{Curriculum learning schedule.}
Activating $\mathcal{L}_\text{RCM}$ from epoch~1 can destabilize early
optimization because the initial predictions are far from
physically consistent solutions.
We therefore adopt a linear warm-up schedule:
\begin{equation}
  \lambda_\text{rcm}(e) =
  \begin{cases}
    0 & e < e_\text{start}, \\
    \lambda_\text{max}\,\dfrac{e - e_\text{start}}{e_\text{end} - e_\text{start}}
      & e_\text{start} \le e \le e_\text{end}, \\
    \lambda_\text{max} & e > e_\text{end},
  \end{cases}
  \label{eq:schedule}
\end{equation}
with $e_\text{start}{=}10$, $e_\text{end}{=}80$, and $\lambda_\text{max}{=}1.0$.
The MSE-only phase in epochs $1$–$10$ drives the network to a reasonable
initial solution, after which the physical constraint progressively tightens
the solution toward RCM consistency.

\subsection{Random Reference-Point Masking}
\label{sec:masking}

At inference time, fewer than three ground reference points may be available.
Training exclusively with all three points visible would create a
training–inference distribution mismatch.
We mitigate this by randomly masking reference points during training:
with probabilities $p_1{=}0.4$, $p_2{=}0.3$, $p_3{=}0.3$,
a training sample retains exactly 1, 2, or 3 reference points, respectively.
Masked channels are zeroed in both the feature vector and the RCM loss weight
mask $w_{b,r}$.
This forces the model to learn trajectory corrections that are robust to
partial RCM observability.

\subsection{Training Configuration}
\label{sec:training}

Training uses the AdamW optimizer \cite{loshchilov2019adamw}
with initial learning rate $\eta{=}10^{-3}$, weight decay $5\times10^{-4}$,
and cosine annealing to $\eta_\text{min}{=}10^{-5}$ over 100 epochs.
The effective batch size is 32 (batch size $8$ per GPU
$\times$ 2 GPUs $\times$ gradient accumulation steps $2$).
Training employs bfloat16 automatic mixed precision (AMP) and
DistributedDataParallel (DDP) across two NVIDIA RTX 4080 SUPER GPUs (32\,GB each).
Early stopping with patience 15 is applied on validation loss.
Sequence subsampling by stride~2 reduces the effective sequence length to
$T{=}2048$ to keep the mLSTM $O(T^2)$ attention matrix within GPU memory bounds.

% ──────────────────────────────────────────────────────────────────────────────
\section{Simulation Dataset}
\label{sec:dataset}

\subsection{Physical Simulation Pipeline}

Each of the $N{=}5{,}000$ training trajectories is generated by a
six-degree-of-freedom (6-DOF) rigid-body simulator of a quadrotor UAV
(mass $m \approx 8$\,kg, $\pm15\%$ random perturbation)
controlled by PID attitude and position loops
(gains randomized $\pm15\%$).
The true APC trajectory $\vp_\text{true}$ is the output of this
closed-loop simulation perturbed by:
(i) sinusoidal wind ($0.05$–$0.25$\,Hz, amplitude $0$–$1.5$\,m/s),
(ii) stochastic gusts (0.2\% probability per step), and
(iii) rigid-body dynamics.
Flight speed is randomized in $[7, 13]$\,m/s and altitude in $[60, 140]$\,m,
yielding total flight durations of $4.096$\,s at PRF $= 1000$\,Hz.

\subsection{Error Injection (v5 Diversity Enhancement)}

The navigation error $\vdelta$ is injected with five components as in
(\ref{eq:error_decomp}).
The overall scale factor $s_\text{err}$ is sampled from a stratified
distribution:
\begin{equation}
  s_\text{err} \sim
  \begin{cases}
    \mathcal{U}(0.2, 0.5) & \text{prob.\ } 0.30 \quad (\text{small}), \\
    \mathcal{U}(0.5, 1.5) & \text{prob.\ } 0.50 \quad (\text{medium}), \\
    \mathcal{U}(1.5, 3.0) & \text{prob.\ } 0.20 \quad (\text{large}).
  \end{cases}
\end{equation}
Step-change fault events ($0$–$5$ per trajectory, magnitude $\pm[20,50]$\,mm)
are superimposed.
Three diversity augmentations prevent the model from memorizing
absolute coordinate statistics:
\begin{enumerate}
  \item \textbf{Random heading rotation} $\theta_h \sim \mathcal{U}(-15^\circ, 15^\circ)$:
        the entire trajectory and reference-point coordinates are rotated in the
        horizontal $(Y_\text{lat}, X_\text{fwd})$ plane.
  \item \textbf{Random lateral offset} $\Delta Y \sim \mathcal{U}(-100, 100)$\,m:
        a constant shift applied to $Y_\text{lat}$ of all points.
  \item \textbf{Random RCM noise} $\sigma_\rho \sim \mathcal{U}(5, 30)$\,mm per sample
        (previously fixed at $3$\,mm).
\end{enumerate}
Reference targets $\{P_A, P_B, P_C\}$ are placed at random positions
spanning the scene and constrained to have pairwise distances $>10$\,m
to avoid rank deficiency in the RCM constraint.
The dataset is split 90\%/10\% into training (4{,}500) and validation (500) sets.

% ──────────────────────────────────────────────────────────────────────────────
\section{Experiments}
\label{sec:experiments}

\subsection{Evaluation Metrics}

The primary metric is the \textbf{position RMSE}:
\begin{equation}
  \text{RMSE} = \sqrt{\frac{1}{NT}\sum_{n,k}
    \norm{\hat{\vp}_{\text{fix},n}[k] - \vp_{\text{true},n}[k]}_2^2},
  \label{eq:rmse}
\end{equation}
evaluated on the validation set.
We also report the \textbf{error reduction rate}
$\eta = 1 - \text{RMSE}_\text{fix}/\text{RMSE}_\text{raw}$.

\subsection{Comparison Methods}

We compare PI-xLSTM against the following baselines,
all trained under identical hyperparameter settings and data conditions:
\begin{itemize}
  \item \textbf{No correction}: Raw navigation output $\vp_\text{raw}$, serving
        as the performance lower bound.
  \item \textbf{Bi-LSTM}: A four-layer bidirectional LSTM \cite{hochreiter1997lstm}
        with the same input features and MSE$+$smooth$+$RCM loss function.
  \item \textbf{Bi-GRU}: A four-layer bidirectional gated recurrent unit \cite{cho2014gru}.
  \item \textbf{Transformer}: A four-layer non-causal Transformer encoder with
        sinusoidal positional encoding, Pre-LN normalization, and
        $H{=}4$ attention heads.
\end{itemize}

\subsection{Ablation Study}

To isolate the contribution of each proposed component, we train three ablated
variants of PI-xLSTM:
\begin{itemize}
  \item \textbf{PI-xLSTM w/o RCM}: $\lambda_\text{max}{=}0$; the physical constraint
        is removed.
  \item \textbf{PI-xLSTM w/o smooth}: $\lambda_s{=}0$; the smoothness regularizer
        is removed.
  \item \textbf{PI-xLSTM w/o RefMask}: The random reference-point masking is
        disabled; all three reference targets are always visible during training.
\end{itemize}

\subsection{Preliminary Results}

\begin{table}[t]
  \centering
  \caption{Validation Trajectory RMSE (preliminary results after 33 epochs).
           Full results will be reported after training convergence
           ($\sim$200 epochs).}
  \label{tab:results}
  \begin{tabular}{lccc}
    \toprule
    Method & RMSE (mm) & $\eta$ (\%) & Params (M) \\
    \midrule
    No correction (P\textsubscript{raw}) & 42.4 & — & — \\
    \midrule
    Bi-LSTM              & \textit{TBD} & \textit{TBD} & 1.3 \\
    Bi-GRU               & \textit{TBD} & \textit{TBD} & 1.0 \\
    Transformer          & \textit{TBD} & \textit{TBD} & 1.5 \\
    \midrule
    PI-xLSTM w/o RCM     & \textit{TBD} & \textit{TBD} & 1.5 \\
    PI-xLSTM w/o Smooth  & \textit{TBD} & \textit{TBD} & 1.5 \\
    PI-xLSTM w/o RefMask & \textit{TBD} & \textit{TBD} & 1.5 \\
    \midrule
    \textbf{PI-xLSTM (proposed)} & \textbf{28.6} & \textbf{32.5} & 1.5 \\
    \bottomrule
  \end{tabular}
\end{table}

Table~\ref{tab:results} shows preliminary validation RMSE
for PI-xLSTM at epoch 33 (training is ongoing toward 200 epochs).
The initial navigation RMSE of 42.4\,mm is reduced to 28.6\,mm,
a 32.5\% improvement achieved within the first 33 epochs.

Fig.~\ref{fig:curves} presents the training dynamics.
The validation loss follows the training loss closely throughout all epochs,
with no evidence of the severe overfitting observed in earlier dataset versions
that lacked heading rotation and lateral offset augmentation.
The RCM physical constraint activates at epoch~10 (as designed) and
produces an immediate acceleration of convergence, with
validation RMSE dropping from 35.5\,mm (epoch~11) to 28.6\,mm (epoch~33).
The gradient norm shows occasional spikes (epochs~16, 28, 30) coinciding with
RCM weight increases; gradient clipping at norm 5.0 prevents divergence.

\begin{figure}[t]
  \centering
  % Replace with actual figure once training converges
  \fbox{\parbox{0.9\linewidth}{\centering\vspace{1.5cm}
        [Training curves figure — to be inserted]\vspace{1.5cm}}}
  \caption{Training and validation losses (top) and validation position RMSE
           (bottom) over training epochs. The dashed vertical line marks
           RCM constraint activation at epoch~10.}
  \label{fig:curves}
\end{figure}

\subsection{Discussion}

\textbf{Role of the RCM physical constraint.}
The RCM loss (\ref{eq:rcm_loss}) provides supervision signal
complementary to the MSE term:
even when the labeled correction $\vdelta$ is noisy or biased,
the physical constraint anchors the predicted trajectory to
the observed slant ranges.
This is particularly important at inference time when
only partial reference-point observations are available.

\textbf{Curriculum learning.}
Immediate activation of $\mathcal{L}_\text{RCM}$ at epoch~1 was found to
cause gradient instability and slower convergence in preliminary experiments.
The linear warm-up schedule (\ref{eq:schedule}) decouples
the initial MSE-based learning phase from
the physically constrained refinement phase,
yielding more stable training dynamics.

\textbf{Training–inference consistency.}
The random reference-point masking strategy is critical for
deployment scenarios where only one or two reference targets are available.
Without masking, the model learns to exploit all three channels and
performs poorly when channels are absent at test time.
With masking, the model learns a graceful degradation strategy:
with one reference point, the correction is less precise but still
substantially better than the raw navigation output
(quantified in the ablation study, TBD).

\textbf{xLSTM vs.\ RNNs and Transformers.}
The mLSTM cell's matrix state enables richer sequence context than
the scalar state of LSTM/GRU, at the cost of quadratic
$O(T^2)$ memory complexity in the parallel mode.
For $T{=}2{,}048$ this is tractable on 32\,GB GPUs (256\,MB per
head per batch element), while a full 4{,}096-step sequence would
require 1\,GB per batch element and exceed available memory.
Compared with Transformer attention, the xLSTM mLSTM attends to the
full history without positional encoding, which we conjecture to
better model the non-stationary navigation error dynamics.

% ──────────────────────────────────────────────────────────────────────────────
\section{Conclusion}
\label{sec:conclusion}

We have presented PI-xLSTM, a physics-informed bidirectional xLSTM network
for UAV-SAR trajectory refinement.
The proposed method combines a 12-dimensional navigation–RCM feature
representation, a compound loss with curriculum-learned RCM physical
consistency, and random reference-point masking for training–inference robustness.
Preliminary results show a 32.5\% trajectory error reduction over the raw
navigation input within 33 training epochs, with no signs of overfitting
thanks to a diversity-enhanced simulation dataset.
Complete comparison with recurrent and attention-based baselines,
ablation studies, and real-flight data validation are ongoing
and will be reported in the final version.

% ──────────────────────────────────────────────────────────────────────────────
\section*{Acknowledgment}
[Acknowledgment text to be added upon acceptance.]

% ──────────────────────────────────────────────────────────────────────────────
\bibliographystyle{IEEEtran}
\bibliography{references}

\end{document}
